\chapter{连续多项式与形式级数}

在本附录中，F 表示 K 上的分离多赋范空间，\(E_{i}\)（\(1 \leqslant i \leqslant n\)）表示 K 上的赋范空间，而 E 表示各 \(E_{i}\) 的拓扑向量空间乘积。

对于 \(\alpha \in \mathbf{N}^n\) 和 \(1 \leqslant j \leqslant |\alpha|\)，置：

\[
\alpha (j) = \inf \left\{k \in \mathbf {N} \mid 1 \leqslant k \leqslant n \text { and } j > \alpha_ {1} + \dots + \alpha_ {k - 1} \right\}
\]

（因此，序列 \(\alpha(j)\) 是将 1 写 \(\alpha_{1}\) 次，再将 \(1,\ldots,\alpha_{n}\) 写 \(n\) 次得到的。）

将 \(E_{\alpha}\) 记为各 \(E_{\alpha(j)}\)（其中 \(1 \leqslant j \leqslant |\alpha| \)）所成的拓扑向量空间乘积。将 \(\operatorname{Hom}_{\alpha}(E_1, \ldots, E_n; F)\) 记为 \(|\alpha|\)-多线性映射空间，从 \(E_{\alpha}\) 到 F，并将 \(\mathcal{L}_{\alpha}(E_1, \ldots, E_n; F)\) 记为 \(\operatorname{Hom}_{\alpha}(E_1, \ldots, E_n; F)\) 的子空间，由连续多线性映射组成，赋予它在 \(E_{\alpha}\) 的有界子集上一致收敛的拓扑；这是一个分离的多赋范空间，其拓扑可由半范数族 \(\| u \|_{\gamma}\) 定义：对 F 上的连续半范数 \(\gamma\)，记 \(\| u \|_{\gamma}\) 为满足下式的数 \(a \geqslant 0\) 的下确界：

\[
\| u (x _ {1}, \dots , x _ {| \alpha |}) \| _ {\gamma} \leqslant a \| x _ {1} \| \dots \| x _ {| \alpha |} \|
\]

其中每个 \(x_{j}\) 遍历 \(E_{\alpha(j)}\)。

将典范投影 \(p_{i}\) 记为从 E 到 \(\mathbf{E}_{i}\) 的映射。对于 \(\alpha\in\mathbf{N}^{n}\)，若 \(\alpha\neq0\)，将 \(p_{\alpha}\) 记为映射 \(x\mapsto(p_{\alpha(j)}(x))\)，从 E 到 \(\mathbf{E}_{\alpha}\)。

A.1. 从 E 到 F 的映射 \(f\) 称为 E 上取值于 F 的多重次数为 \(\alpha\) 的多齐次多项式（其中 \(\alpha \in \mathbf{N}^{n}\)），如果存在一个元素

\[
u \in \operatorname{Hom} _ {\alpha} (\mathrm{E} _ {1}, \dots , \mathrm{E} _ {n}; \mathrm{F})
\]

使得 \(f = u \circ p_{\alpha}\)。

A.2. 将 \(\mathrm{P}_{\alpha}(\mathrm{E}_{1},\ldots,\mathrm{E}_{n};\mathrm{F})\) 记为 \(\mathcal{L}_{\alpha}(\mathrm{E}_{1},\ldots,\mathrm{E}_{n};\mathrm{F})\) 在映射 \(u\mapsto u\circ p_{\alpha}\)（它是从 \(\mathrm{Hom}_{\alpha}(\mathrm{E}_{1},\ldots,\mathrm{E}_{n};\mathrm{F})\) 到从 E 到 F 的映射空间的线性映射）下的像，并赋予它由 \(\mathcal{L}_{\alpha}(\mathrm{E}_{1},\ldots,\mathrm{E}_{n};\mathrm{F})\) 的拓扑诱导的商拓扑。\(\mathrm{P}_{\alpha}(\mathrm{E}_{1},\ldots,\mathrm{E}_{n};\mathrm{F})\) 的元素称为 E 上取值于 F 的多重次数为 \(\alpha\) 的多齐次连续多项式。\(\mathrm{P}_{\alpha}(\mathrm{E}_{1},\ldots,\mathrm{E}_{n};\mathrm{F})\) 的拓扑由以下半范数族定义：

\[
\| f \| _ {\gamma} = \inf _ {u \in \mathscr {L} _ {\alpha} (\mathrm{E} _ {1}, \dots , \mathrm{E} _ {n}; \mathrm{F}), f = u \circ p _ {\alpha}} \| u \| _ {\gamma}
\]

其中 \(\gamma\) 遍历 F 上的连续半范数集。若 F 是赋范空间，范数为 \(\gamma\)，则用 \(\|f\|\) 代替 \(\|f\|_{\gamma}\)。空间 \(P_{\alpha}(E_{1},\ldots,E_{n};F)\) 及其拓扑不因将各 \(E_{i}\) 上给定的范数换成等价范数而改变。因此，当各 \(E_{i}\) 是可赋范的拓扑向量空间时，也可以定义它们。

特别地，\(\mathbf{P}_{k}(\mathbf{E};\mathbf{F})\) 的元素称为总次数为 \(k\) 的齐次连续多项式，在 \(\mathbf{E}\) 上取值于 \(\mathbf{F}\)。空间 \(\mathbf{P}_{k}(\mathbf{E};\mathbf{F})\) 是各空间 \(\mathbf{P}_{\alpha}(\mathbf{E}_{1},\ldots,\mathbf{E}_{n};\mathbf{F})\) 的拓扑直和，其中 \(|\alpha| = k\)。空间 \(\mathbf{P}_{0}(\mathbf{E};\mathbf{F})\) 是从 \(\mathbf{E}\) 到 \(\mathbf{F}\) 的常值映射空间，并与 \(\mathbf{F}\) 认同。

A.3. 将 P(E; F) 或 P(E₁, \ldots, Eₙ; F) 记为从 E 到 F 的所有映射组成的向量空间中由子空间 Pₖ(E; F) 生成的向量子空间。

它既是各子空间 \(\mathbf{P}_{\alpha}(\mathbf{E}_{1},\ldots,\mathbf{E}_{n};\mathbf{F})\) 的直和，其中 \(\alpha\in\mathbf{N}^{n}\)，也是各子空间 \(\mathbf{P}_{k}(\mathbf{E};\mathbf{F})\) 的直和，其中 \(k\in\mathbf{N}\)。\(\mathbf{P}(\mathbf{E};\mathbf{F})\) 的元素称为 \(\mathbf{E}\) 上取值于 \(\mathbf{F}\) 的连续多项式。

A.4. 设 \(G_{j}\) 是赋范空间（\(1 \leqslant j \leqslant m\)）。设 \(f_{j} \in \mathrm{P}(\mathrm{E}_{1}, \ldots, \mathrm{E}_{n}; \mathrm{G}_{j})\)（\(1 \leqslant j \leqslant m\)），且 \(g \in \mathrm{P}(\mathrm{G}_{1}, \ldots, \mathrm{G}_{m}; \mathrm{F})\)。映射

\[
h: x \mapsto g (f _ {1} (x), \dots , f _ {m} (x))
\]

属于 \(\mathbf{P}(\mathbf{E}_1,\ldots ,\mathbf{E}_n;\mathbf{F})\)。此外，若 \(f_{j}\in \mathbf{P}_{\alpha}(\mathbf{E}_{1},\ldots ,\mathbf{E}_{n};\mathbf{G}_{j})\)，对 \(1\leqslant j\leqslant m\)，且 \(\alpha \in \mathbf{N}^n\)，又有 \(g\in \mathbf{P}_{\beta}(\mathbf{G}_1,\ldots ,\mathbf{G}_m;\mathbf{F})\)，其中 \(\beta \in \mathbf{N}^{m}\)，则 \(h\in \mathbf{P}_{|\beta |\alpha}(\mathbf{E}_1,\ldots ,\mathbf{E}_n;\mathbf{F})\)，并且

\[
\| h \| _ {\gamma} \leqslant \| g \| _ {\gamma} \cdot \| f \| ^ {\beta}
\]

对 F 上的每个连续半范数 \(\gamma\) 都成立（置 \(\| f\|^{\beta} = \prod_{1 \leqslant j \leqslant m} \| f_j \|^{\beta_j}\)）。

A.5. 将 \(\hat{\mathbf{P}}(\mathbf{E}_{1},\ldots,\mathbf{E}_{n};\mathbf{F})\) 记为各 \(\mathbf{P}_{\alpha}(\mathbf{E}_{1},\ldots;\mathbf{E}_{n};\mathbf{F})\) 的乘积向量空间，其中 \(\alpha\in\mathbb{N}^{n}\)，并赋予各因子离散拓扑的乘积拓扑。这个拓扑使 \(\hat{\mathbf{P}}(\mathbf{E}_{1},\ldots,\mathbf{E}_{n};\mathbf{F})\) 成为分离且完备的拓扑群。从 \(\mathbf{P}(\mathbf{E}_{1},\ldots,\mathbf{E}_{n};\mathbf{F})\) 到 \(\hat{\mathbf{P}}(\mathbf{E}_{1},\ldots,\mathbf{E}_{n};\mathbf{F})\) 的线性映射延拓了各 \(\mathbf{P}_{\alpha}(\mathbf{E}_{1},\ldots,\mathbf{E}_{n};\mathbf{F})\) 的典范嵌入；它是单射，且其像在 \(\hat{\mathbf{P}}(\mathbf{E}_{1},\ldots,\mathbf{E}_{n};\mathbf{F})\) 中稠密：通常将 E 上取值于 F 的连续多项式与其在 \(\hat{\mathbf{P}}(\mathbf{E}_{1},\ldots,\mathbf{E}_{n};\mathbf{F})\) 中的像认同。

\(\mathbf{P}(\mathbf{E}_{1},\ldots,\mathbf{E}_{n};\mathbf{F})=\mathbf{P}(\mathbf{E};\mathbf{F})\) 的恒等映射按连续性延拓为从 \(\hat{\mathbf{P}}(\mathbf{E}_{1},\ldots,\mathbf{E}_{n};\mathbf{F})\) 到 \(\hat{\mathbf{P}}(\mathbf{E};\mathbf{F})\) 的同构。

\(\hat{\mathbf{P}}(\mathbf{E}_{1},\ldots,\mathbf{E}_{n};\mathbf{F})\) 的元素称为在各 \(\mathbf{E}_{i}\) 的乘积上取值于 \(\mathbf{F}\) 的具有连续分量的形式级数（或语言上滥称为形式级数）。

A.6. 设 \(\mathbf{G}_j\)（\(1 \leqslant j \leqslant m\)）是赋范空间，且 \(g \in \mathrm{P}(\mathrm{G}_1, \ldots, \mathrm{G}_m; \mathrm{F})\)。映射 \(\mathbf{f} \mapsto g \circ \mathbf{f}\)，从 \(\prod_{1 \leqslant j \leqslant m} \mathrm{P}(\mathrm{E}_1, \ldots, \mathrm{E}_n; \mathrm{G}_j)\) 到 \(\mathrm{P}(\mathrm{E}_1, \ldots, \mathrm{E}_n; \mathrm{F})\)，按连续性延拓为映射 \(\mathbf{f} \mapsto g \circ \mathbf{f}\)，从 \(\prod_{1 \leqslant j \leqslant m} \hat{\mathrm{P}}(\mathrm{E}_1, \ldots, \mathrm{E}_n; \mathrm{G}_j)\) 到 \(\hat{\mathrm{P}}(\mathrm{E}_1, \ldots, \mathrm{E}_n; \mathrm{F})\)。

设 \(f_j = (f_{j,\alpha})_{\alpha \in \mathbb{N}^n} \in \hat{\mathbf{P}}(\mathbf{E}_1, \ldots, \mathbf{E}_n; \mathbf{G}_j)\)，且 \(f_{j,0} = 0\)。置 \(\mathbf{f} = (f_j)\)：映射 \(g \mapsto g \circ \mathbf{f}\)，从 \(\mathrm{P}(\mathbf{G}_1, \ldots, \mathbf{G}_m; \mathbf{F})\) 到 \(\hat{\mathbf{P}}(\mathbf{E}_1, \ldots, \mathbf{E}_n; \mathbf{F})\)，按连续性延拓为映射 \(g \mapsto g \circ \mathbf{f}\)，从 \(\hat{\mathbf{P}}(\mathbf{G}_1, \ldots, \mathbf{G}_m; \mathbf{F})\) 到 \(\hat{\mathbf{P}}(\mathbf{E}_1, \ldots, \mathbf{E}_n; \mathbf{F})\)。

